\documentclass[12pt]{article}

\usepackage[spanish]{babel}
\usepackage[utf8]{inputenc}
\usepackage[T1]{fontenc}
\usepackage{geometry}
\usepackage{graphicx}
\usepackage{setspace}
\usepackage{titlesec}
\usepackage{listings}
\usepackage{xcolor}
\usepackage{float}
\usepackage{hyperref}

\geometry{margin=1in}
\onehalfspacing

\definecolor{codegray}{rgb}{0.95,0.95,0.95}

\lstset{
    backgroundcolor=\color{codegray},
    basicstyle=\ttfamily\small,
    breaklines=true,
    frame=single,
    language=Scala
}

\begin{document}

% portads

\begin{titlepage}
    \centering

    % Logo opcional
    %\includegraphics[width=0.3\textwidth]{logo.png}\par\vspace{1cm}

    {\scshape\Large Programación Funcional \par}
    \vspace{1.5cm}

    {\huge\bfseries Evidencia del Trabajo \par}
    \vspace{0.5cm}

    {\Large Proyecto Nones: Simulación Bioclimática \par}

    \vfill

    {\large
    Nombre del Alumno: Emilio Izquierdo Montero \\
    \vspace{0.3cm}

    Profesor: Octavo Guzman Peñaloza \\
    \vspace{0.3cm}

    Grupo / Semestre: 8 \\
    \vspace{0.3cm}

    Fecha: 10 de mayo del 2026
    \par}

    \vfill

\end{titlepage}


\section{Introducción}

El presente proyecto implementa una simulación bioclimática utilizando el lenguaje Scala y principios de programación funcional, el sistema genera datos simulados provenientes de sensores ambientales y clasifica el estado de un ecosistema mediante un arbol de decisión funcional.

El desarrollo del proyecto hace uso de estructuras inmutables, evaluación perezosa mediante \texttt{LazyList} y procesamiento funcional utilizando transformaciones como \texttt{map} y recursividad.


\section{Objetivo}

Desarrollar un sistema funcional capaz de:

\begin{itemize}
    \item Simular datos de sensores ambientales.
    \item Procesar información usando programación funcional.
    \item Implementar un árbol de decisión para clasificar estados.
    \item Utilizar Lazy Evaluation con \texttt{LazyList}.
\end{itemize}


\section{Modelado de Datos}

Se utilizaron ADTs (Algebraic Data Types) para representar tanto los datos de sensores como los estados del ecosistema.

\begin{lstlisting}
case class SensorData(
    id: Int,
    temp: Double,
    humedad: Double,
    co2: Int
)

sealed trait EstadoSalud

case object Estable extends EstadoSalud
case object Alerta extends EstadoSalud
case object Critico extends EstadoSalud
\end{lstlisting}


\section{Árbol de Decisión}

El sistema clasifica el estado del ecosistema mediante un árbol de decisión funcional.

Las reglas implementadas fueron:

\begin{itemize}
    \item CO2 mayor a 800 ppm = Crítico
    \item Temperatura mayor a 32°C = Alerta
    \item Humedad menor a 25\% = Alerta
    \item Cualquier otro caso = Estable
\end{itemize}

\begin{lstlisting}
sealed trait DecisionTree

case class DecisionNode(
    condicion: SensorData => Boolean,
    izquierda: DecisionTree,
    derecha: DecisionTree
) extends DecisionTree

case class Hoja(resultado: EstadoSalud)
extends DecisionTree
\end{lstlisting}


\section{Uso de LazyList}

Se utilizó \texttt{LazyList} para generar un flujo infinito de sensores simulados.

Esto permite que los datos sean calculados únicamente cuando son necesarios, aplicando una evaluación bastanteeee lazy.

\begin{lstlisting}
def streamSensores: LazyList[SensorData]
\end{lstlisting}


\section{Evidencias de Ejecución}

\subsection{Compilación y ejecución del proyecto}

A continuacion se mostraran capturas de:

\begin{itemize}
    \item La terminal con sbt instalado.
    \item El comando \texttt{sbt run}
    \item La compilación correcta
\end{itemize}

\vspace{0.5cm}

\begin{figure}[H]
    \centering
    \includegraphics[width=0.9\textwidth]{/home/emilio/Descargas/scalafinal/creacionsbt.png}
\end{figure}

\vspace{0.5cm}

\begin{figure}[H]
    \centering
    \includegraphics[width=0.9\textwidth]{/home/emilio/Descargas/scalafinal/sbtrun.png}
\end{figure}


\subsection{Salida del programa}

Las imagenes que siguen es donde se observa:

\begin{itemize}
    \item Los sensores generados
    \item Temperatura
    \item Humedad
    \item CO2
    \item Estado clasificado
\end{itemize}

\begin{figure}[H]
    \centering
    \includegraphics[width=0.9\textwidth]{/home/emilio/Descargas/scalafinal/estados.png}
\end{figure}

\section*{Código Final}

El código final desarrollado en nvim usando sbt, metals se encuentra disponible en el repositorio de GitHub:

\begin{itemize}
    \item Repositorio: \href{https://github.com/emili0p/funcprog/tree/master/proyectonones}{https://github.com/emili0p/funcprog/tree/master/proyectonones}
\end{itemize}

\noindent
Se incluyen archivos \texttt{.scala} correspondientes al proyecto, que pueden descargarse y ejecutarse directamente en donde sea usando sbt run para reproducir los resultados mostrados en las capturas de pantalla.

\section{Conclusión}

Scala permite implementar sistemas funcionales de forma clara y modular. El uso de ADTs, recursividad y \texttt{LazyList} facilitó la creación de un sistema eficiente para simular datos ambientales y clasificarlos mediante un árbol de decisión funcional.


\end{document}
